 % ---------------------------------------------------------------------
%  Chapitre 3 — BLOC 1 (chapitre développé)
%  Suit les 5 points des consignes pour CHAQUE méthode :
%   théorie -> conditions d'application -> résultats -> limites
% ---------------------------------------------------------------------
\chapter{Bloc 1 --- valeur des stratégies techniques}
\label{chap:bloc1}

Ce chapitre traite la question~: \emph{parmi dix stratégies d'analyse
technique, lesquelles produisent un avantage (\edgevar) significativement
positif~?} L'analyse procède en six étapes, des tests univariés aux
analyses factorielles.

% =====================================================================
\section{Étape 1 --- chaque stratégie bat-elle le hasard ?}
\label{sec:b1-etape1}

\subsection{La question posée à l'outil}
% [À RÉDIGER — schéma « on se demande… »]
% Idées à développer avec TES mots :
%  - Première brique du protocole de validation : une stratégie qui
%    affiche un edge moyen positif le doit-elle à un vrai avantage, ou
%    à la variabilité d'échantillonnage ?
%  - Formaliser : comparer la moyenne de l'edge à la valeur 0.
%    H0 : mu_edge = 0 (pas mieux que le hasard) ; H1 : mu_edge > 0.
%  - Pourquoi unilatéral (> 0) : on ne s'intéresse qu'à un avantage,
%    pas à un désavantage.
Pour évaluer une stratégie, on commence par la comparer au hasard. À
chaque trade réalisé par la stratégie, on associe 200~trades fictifs de
même exposition (même titre, même durée, même sens), mais dont la date
d'entrée est tirée au sort~; la moyenne de leurs rendements définit le
\emph{gain du hasard} pour ce trade.

On définit alors, pour chaque trade, son \edgevar{} comme le rendement
effectivement obtenu par la stratégie (éventuellement négatif) diminué de
ce gain du hasard~:
\[
  \edgevar = \text{gain de la stratégie} - \text{gain du hasard}.
\]
Un \edgevar{} positif signifie que la stratégie a fait mieux qu'une entrée
au hasard de même exposition~: son \emph{timing} apporte une information
que la seule dérive du marché n'explique pas.

La première question posée à l'outil est donc~: l'avantage moyen affiché
par une stratégie est-il réel, ou n'est-il qu'un accident
d'échantillonnage~? On la formalise en confrontant la moyenne de
l'\edgevar{} à la valeur zéro, ce qui conduit au test de Student présenté
ci-dessous.

\subsection{Les méthodes qui y répondent}

% --- MÉTHODE 1 : le test principal -----------------------------------
\subsubsection*{Test de Student}

Lorsque l'on obtient un edge moyen positif, on cherche à savoir si la
stratégie constitue réellement une prédiction, ou si cela est dû à l'aléa
de l'échantillon.

Il faut ici distinguer deux quantités~: l'edge moyen
\emph{théorique} \(\mu_{\edgevar}\), c'est-à-dire la vraie valeur (inconnue)
gouvernant la stratégie, et sa moyenne
\emph{observée} \(\bar e\), calculée sur notre échantillon de trades
(connue). On se pose la question suivante~: à partir de la valeur observée
\(\bar e\), peut-on conclure que la vraie valeur \(\mu_{\edgevar}\) est
positive~?

Le test de Student à un échantillon répond à cette question en confrontant
\(\mu_{\edgevar}\) à une valeur de référence fixée à zéro (zéro~$=$ aucun
avantage sur le hasard). Sa mise en \oe uvre suppose toutefois que l'edge
vérifie un critère de normalité~; cette condition, d'abord énoncée sur
l'edge lui-même, sera ensuite assouplie grâce au théorème central limite,
comme on le verra à la suite de ce test. On oppose ainsi deux
hypothèses~:
\[
  \begin{cases}
    \Hzero : \mu_{\edgevar} = 0 & \text{(le vrai edge est nul : pas mieux que le hasard),}\\[2pt]
    \Hun : \mu_{\edgevar} > 0 & \text{(le vrai edge est positif : avantage réel).}
  \end{cases}
\]
Le test est \emph{unilatéral}~: seul un avantage strictement positif nous
intéresse.

Pour trancher, on définit \(t\) comme étant le nombre d'erreurs-types
nécessaires pour faire basculer l'edge moyen de zéro vers celui obtenu
dans notre échantillon~:
\[
  t = \frac{\bar e}{s/\sqrt{n}} \sim t_{n-1} \quad\text{sous } \Hzero .
\]

Sous \Hzero{} (edge nul), \(t\) suit une loi de Student à \(n-1\) degrés
de liberté. On évalue alors la p-value, c'est-à-dire la probabilité
d'obtenir un \(t\) au moins aussi grand que celui observé si \Hzero{}
était vraie.

Si la p-value est très faible, le hasard n'est plus une
explication crédible de l'edge moyen positif observé.

On tranche alors sur le fait que l'edge moyen théorique est positif, et que la
stratégie permet donc de faire mieux que le hasard.


On a dit précédemment que le test de Student nécessitait que l'edge suive
une loi normale. On cherche à vérifier cette condition, dans un premier
temps avec le test de Shapiro--Wilk, puis dans un second temps à l'aide du
théorème central limite (TCL). En effet, le test de Shapiro--Wilk a peu de
chances de valider cette hypothèse~: la taille conséquente de l'échantillon
déstabilise ce test, alors qu'elle renforce au contraire le TCL.

% --- MÉTHODE 2 : la condition du test principal ----------------------
\subsubsection*{Test de normalité de Shapiro--Wilk} % [À RÉDIGER]
% Idées à mettre dans TES mots :
%  - Enchaînement : le Student suppose (en toute rigueur) des données
%    normales ; avant de l'appliquer, l'outil doit donc TESTER cette
%    hypothèse.
%  - Shapiro compare une estimation de la variance fondée sur les
%    statistiques d'ordre à l'estimation classique ; W proche de 1 =
%    échantillon compatible avec une loi normale.
Le test de Shapiro--Wilk est un test de normalité qui consiste à comparer
deux mesures de dispersion~: d'une part les écarts à la moyenne, d'autre
part l'ordre des valeurs. Ces deux mesures ne coïncident que si
l'hypothèse de normalité est vraie. On oppose ainsi deux hypothèses~:
\[
  \begin{cases}
    \Hzero \;(W \simeq 1) : \text{l'edge suit une loi normale,}\\[2pt]
    \Hun \;(W \ll 1) : \text{l'edge ne suit pas une loi normale.}
  \end{cases}
\]
On notera que, contrairement au test de Student, c'est ici \Hzero{} qui
correspond à la propriété recherchée (la normalité)~: une p-value faible
conduira donc à \emph{rejeter} la normalité.

On pose alors la statistique \(W = b^2 / \mathrm{SCE}\) comme rapport de
ces deux mesures. Lorsque \(W\) est proche de 1, l'échantillon est
compatible avec une loi normale.

Pour trancher sur la normalité, on procède de façon similaire au test de
Student~: on calcule la p-value d'obtenir notre \(W\) observé si la loi
était normale~; si cette probabilité est trop faible, alors on rejette
l'hypothèse de normalité.

% Constat sur nos données (détail W + p-value en annexe)
Appliqué à nos données, le test rejette la normalité pour les dix
stratégies (statistique \(W\) et p-values détaillées en
annexe~\ref{tab:ann-shapiro}). La condition du test de Student n'est donc,
en apparence, pas remplie.

% --- LE TCL : ce qui débloque le test de Student ---------------------
\subsubsection*{Théorème central limite} % [À RÉDIGER]
% Idées à mettre dans TES mots :
%  - Le rejet de Shapiro porte sur l'edge INDIVIDUEL (non normal :
%    queues épaisses des rendements financiers).
%  - MAIS le test de Student ne dépend que de la MOYENNE (edge-bar),
%    pas des edges individuels.
%  - Le TCL garantit que la moyenne d'un grand échantillon est
%    quasi-normale QUELLE QUE SOIT la loi de départ.
%  - Donc : Student reste valide malgré le rejet de Shapiro, parce que
%    c'est la normalité de la MOYENNE qui compte, et elle est assurée
%    par le grand n (973 à 8469 trades).
%  - Bien insister : le TCL ne dit PAS "les données sont normales",
%    il dit "la moyenne des données est normale".
Le rejet de Shapiro--Wilk porte sur la non-normalité de l'edge
individuel. Mais comme le test de Student ne dépend que de la
moyenne de l'edge, on peut se rabattre sur le théorème central
limite, qui garantit la quasi-normalité de la moyenne sur les grands
échantillons, quelle que soit la loi de départ~:
\[
  \frac{\bar e - \mu_{\edgevar}}{\sigma/\sqrt{n}}
  \;\xrightarrow[n \to \infty]{}\; \mathcal{N}(0,1) .
\]
La condition réellement utile au test de Student est donc satisfaite,
malgré le rejet de Shapiro--Wilk.

On remarque de plus que cette quantité est exactement celle du test de
Student~: sous \Hzero{} (\(\mu_{\edgevar} = 0\)), et en remplaçant l'écart-type
théorique \(\sigma\) (inconnu) par son estimation \(s\), on retrouve la
statistique \(t = \bar e / (s/\sqrt{n})\). Le théorème central limite
justifie ainsi directement la forme et la validité du test de Student.

% --- MÉTHODE 3 : le garde-fou des tests multiples --------------------
\subsubsection*{Correction de Bonferroni} % [À RÉDIGER]
% Idées à mettre dans TES mots :
%  - Enchaînement : comme l'outil teste plusieurs stratégies d'un coup,
%    il faut se protéger des faux positifs.
%  - Tester 10 stratégies au seuil 5 %, c'est multiplier les occasions
%    d'un faux positif (P(au moins un) ~ 40 %). Bonferroni durcit le
%    seuil individuel pour contrôler le risque GLOBAL (inégalité de Boole).
Pour appliquer le test de Student, il faut fixer une valeur seuil pour la
p-value. Lorsque l'on effectue un test isolé, on prend habituellement
\(0{,}05\)~: on estime que s'il n'y avait que \(5\,\%\) de chances
d'obtenir un tel résultat sous \Hzero{}, il est peu vraisemblable d'être
tombé dessus par hasard, et l'on accepte alors le risque de rejeter
\Hzero{}.

Seulement, ici, on effectue \(10\) tests. Plaçons-nous dans le cas le plus
défavorable où aucune stratégie n'a d'avantage réel (\Hzero{} vraie pour
les dix)~: pour un test isolé, la probabilité de ne pas conclure à tort à
un avantage est alors \(0{,}95\). En supposant les tests indépendants, la
probabilité de ne commettre aucune erreur sur les \(10\) tests tombe à
\(0{,}95^{10} \approx 0{,}60\). Le risque de commettre \emph{au moins un}
faux positif sur l'ensemble s'élève donc à
\(1 - 0{,}95^{10} \approx 0{,}40\), ce qui est inacceptable.

Pour éviter cela, on raisonne en sens inverse~: on fixe le risque
\emph{global} à \(0{,}05\) et l'on en déduit le seuil de chaque test.
On s'appuie sur l'\textbf{inégalité de Boole}, selon laquelle la
probabilité qu'au moins un événement se produise est inférieure ou égale
à la somme de leurs probabilités individuelles. Pour garantir un risque
global inférieur à \(0{,}05\), il suffit donc que la somme des seuils
individuels soit égale à \(0{,}05\). En les prenant tous égaux, on obtient
la \textbf{correction de Bonferroni}~:
\[
  \alpha' = \frac{\alpha}{m} = \frac{0{,}05}{10} = 0{,}005 .
\]

\begin{conditions}
Indépendance des trades~; normalité \emph{de la moyenne} assurée par le
théorème central limite ($n$ compris entre 973 et 8\,469 selon la
stratégie).
\end{conditions}

\subsection{Résultats}
% Chiffres RÉELS extraits de bloc1/03_resultats/etape1_tests.txt
\begin{table}[H]
\centering
\caption{Test de l'avantage moyen par stratégie (seuil Bonferroni $=0{,}005$).}
\label{tab:b1-etape1}
\small
\begin{tabular}{@{}lrrrc@{}}
\toprule
Stratégie & $n$ & \edgevar{} moyen & $p$ (Student) & Verdict \\
\midrule
\texttt{rsi\_classic} & 7\,083 & $+0{,}0127$ & $3{,}9\times10^{-9}$  & \textbf{bat le hasard} \\
\texttt{rsi\_strict}  &   973  & $+0{,}1308$ & $8\times10^{-10}$     & \textbf{bat le hasard} \\
\texttt{rsi\_trend}   & 2\,777 & $+0{,}0027$ & $0{,}19$             & non \\
\texttt{db\_bottom}   & 8\,469 & $-0{,}0029$ & $0{,}99$             & non \\
\texttt{ma\_crossover}& 5\,411 & $-0{,}0241$ & $\approx 1$          & non \\
\multicolumn{5}{c}{\footnotesize (autres stratégies~: $p>0{,}45$, non significatives)} \\
\bottomrule
\end{tabular}
\end{table}

Seules \texttt{rsi\_classic} et \texttt{rsi\_strict} survivent à la
correction de Bonferroni.

\begin{limites}
À très grand $n$, un test de Student devient sensible au moindre écart~:
une p-value très faible signale un edge \emph{statistiquement} non nul,
sans garantir qu'il soit \emph{économiquement} exploitable (frais,
exécution). L'ampleur de l'edge, et non sa seule significativité, devra
donc être prise en compte.
\end{limites}

% =====================================================================
\section{Étape 2 --- les stratégies diffèrent-elles entre elles ?}
\label{sec:b1-etape2}

\subsection{Présentation théorique}
% [À DÉVELOPPER] ANOVA à 1 facteur : décomposition SCE_inter / SCE_intra,
% statistique F ~ Fisher(k-1, n-k). Puis Tukey HSD (range studentisé).
% Puis ANOVA à 2 facteurs (strategy x regime) avec terme d'interaction.
\paragraph{ANOVA à un facteur.}
\[
  F = \frac{\mathrm{SCE}_{\text{inter}}/(k-1)}
           {\mathrm{SCE}_{\text{intra}}/(n-k)} \sim F_{k-1,\,n-k}.
\]
La décomposition de la variance et le test \(F\) suivent le cadre
classique de l'analyse de la variance~\cite{scheffe1959}.

\paragraph{Comparaisons multiples (Tukey HSD).}
Fondées sur la loi du \emph{range studentisé}, elles corrigent
automatiquement le risque lié aux $k(k-1)/2$ paires.

\paragraph{ANOVA à deux facteurs avec interaction.}
Modèle \(\edgevar = \mu + \alpha_i + \beta_j + \gamma_{ij} + \varepsilon\)~;
le terme d'interaction \(\gamma_{ij}\) teste si l'effet de la stratégie
\emph{dépend du régime de marché}.

\begin{conditions}
Indépendance~; normalité des résidus (TCL)~; \textbf{homoscédasticité}
vérifiée par les tests de \textbf{Bartlett} et de \textbf{Levene}. En cas
de variances inégales, bascule possible vers l'ANOVA de Welch.
\end{conditions}

\subsection{Résultats}
% Chiffres réels : F ~ 47, p ~ 1e-85 (1 facteur) ; interaction p ~ 5e-69
L'ANOVA à un facteur rejette massivement l'égalité des moyennes
(\(F\approx 47\), \(p\approx 10^{-85}\)). Le test de Tukey identifie
18~paires significatives sur 45, \texttt{rsi\_strict} se détachant du
reste. L'interaction stratégie~$\times$~régime est très significative
(\(p\approx 5\times10^{-69}\))~: \emph{la valeur d'une stratégie dépend
du régime de marché}.

% =====================================================================
\section{Étape 3 --- gagner dépend-il du contexte ?}
\label{sec:b1-etape3}

\subsection{Présentation théorique}
% [À DÉVELOPPER] Test du khi-deux d'indépendance sur tableau de
% contingence : effectifs attendus E_ij = (N_i. N_.j)/N ;
% chi2 = somme (N_ij - E_ij)^2 / E_ij ~ chi2 à (r-1)(c-1) ddl.
\[
  \chi^2 = \sum_{i,j}\frac{(N_{ij}-E_{ij})^2}{E_{ij}},
  \qquad E_{ij}=\frac{N_{i\cdot}N_{\cdot j}}{N},
  \qquad \chi^2 \sim \chi^2_{(r-1)(c-1)}.
\]

\begin{conditions}
Effectifs théoriques \(E_{ij}\ge 5\) (règle de Cochran). Sur tableaux
\(2\times2\) ou à faibles effectifs, recours possible à la correction de
\textbf{Yates}, au \textbf{test exact de Fisher} ou au \textbf{G-test}.
\end{conditions}

\subsection{Résultats}
% Chiffres réels de etape3_chi2.txt
\begin{table}[H]
\centering
\caption{Dépendance de \texttt{win} au contexte (khi-deux d'indépendance).}
\label{tab:b1-etape3}
\small
\begin{tabular}{@{}lrrc@{}}
\toprule
Croisement & $\chi^2$ & ddl & Verdict \\
\midrule
\texttt{win} $\times$ \texttt{regime\_entry} & $112{,}3$   & 2  & lien significatif \\
\texttt{win} $\times$ \texttt{sector}        & $16{,}6$    & 10 & indépendance ($p=0{,}08$) \\
\texttt{win} $\times$ \texttt{strategy}      & $6\,381{,}7$& 9  & lien très fort \\
\bottomrule
\end{tabular}
\end{table}

Le taux de gain dépend fortement de la \emph{stratégie}, faiblement du
\emph{régime}, et pas du \emph{secteur}.

\begin{limites}
À très grand \(n\), le \(\chi^2\) devient significatif même pour des
liens négligeables~; la \emph{significativité} ne dit rien de la
\emph{force} du lien. % justifie la prudence d'interprétation
\end{limites}

% =====================================================================
\section{Étape 4 --- ACP du contexte d'entrée}
\label{sec:b1-etape4}

\subsection{Présentation théorique}
% [À DÉVELOPPER] ACP normée (Périnel) : centrage-réduction,
% diagonalisation de la matrice de corrélation R, valeurs propres
% lambda_k = inertie, cercle des corrélations, cos^2, CTR.
% Critères Kaiser / Cattell pour le nombre d'axes. Variables
% illustratives (régime, secteur, win) projetées a posteriori (V.test).
On réalise une \textbf{ACP normée}~\cite{saporta2006,lebart2006} sur les
variables quantitatives de contexte (\texttt{vol\_entry},
\texttt{rsi\_entry}, \texttt{dist\_ma200},
\texttt{holding\_days})~; le nombre d'axes est justifié par les critères
de \textbf{Kaiser} et de \textbf{Cattell}. Les variables qualitatives
(\texttt{regime}, \texttt{secteur}, \texttt{win}) sont projetées en
\textbf{illustratives} (valeur-test).

\begin{conditions}
Variables quantitatives~; standardisation obligatoire (unités
hétérogènes)~; l'ACP ne capture que des relations \emph{linéaires}.
\end{conditions}

\subsection{Résultats}
% [À COMPLÉTER avec etape4_acp.txt / etape4_acp.png]
\emph{[Insérer la figure \texttt{etape4\_acp.png}, l'inertie des axes et
la lecture du plan factoriel.]}

% =====================================================================
\section{Étapes 5 et 6 --- AFC et ACM}
\label{sec:b1-etape56}

\subsection{Présentation théorique}
% [À DÉVELOPPER] AFC sur tableau de contingence strategy x tranche de
% résultat (profils lignes/colonnes, distance du chi-deux, résidus de
% Haberman). ACM sur plusieurs variables qualitatives (correction de
% Benzécri, rapport de corrélation eta^2).
L'\textbf{AFC} analyse le tableau de contingence croisant stratégie et
tranche de résultat~; l'\textbf{ACM} étend l'analyse à plusieurs
variables qualitatives simultanément.

\subsection{Résultats}
% [À COMPLÉTER avec etape5_afc.* et etape6_acm.*]
\emph{[Insérer les plans factoriels AFC et ACM et leur interprétation.]}

% =====================================================================
\section{Synthèse du Bloc 1}
\label{sec:b1-synthese}
Sur dix stratégies, seules les deux variantes RSI (hors filtre de
tendance) démontrent un avantage réel~; cet avantage dépend du régime de
marché. Le secteur n'influe pas sur la probabilité de gain.
